\documentclass[12pt]{article}
\usepackage[margin=2cm]{geometry}

\usepackage{graphicx}
\usepackage{algorithm}
\usepackage{algpseudocode}
\usepackage{amsmath}

\usepackage{xeCJK}
\setmainfont{Times New Roman}
\setCJKmainfont[AutoFakeBold=true,AutoFakeSlant=true]{TW-Kai}
\renewcommand\contentsname{目錄}
\renewcommand\listfigurename{圖目錄}
\renewcommand\listtablename{表目錄}
\renewcommand\abstractname{摘要}
\renewcommand\figurename{圖}
\renewcommand\tablename{表}
\renewcommand\refname{參考文獻}

\usepackage{indentfirst}
\setlength{\parindent}{0em}
\setlength{\parskip}{0.5em}

\usepackage[hidelinks]{hyperref}
\newcommand*{\fullref}[1]{\hyperref[{#1}]{\ref*{#1} \nameref*{#1}}}

\renewcommand{\figureautorefname}{圖}
\renewcommand\tableautorefname{表}

\begin{document}

\title{Physical Design for Nanometer ICs\\Programming Assignment \#1\\2-Way F-M Circuit Partitioning}
\author{臺科大 B11107051 李品翰}
\date{2026-03-31}
\maketitle

\section{Partition 主要程式}

本階段負責讀取測資、初始化資料，並作為進入多層次分割 (Multilevel Partitioning) 的入口。

\begin{algorithm}[H]
\caption{Main Partition Flow}
\begin{algorithmic}[1]
\State 初始化所有 cell 的 size 為 1
\State 呼叫 \text{hMetis} 進行多層次分割
\State 取得最終的 cut size 與 cell 分組結果並輸出
\end{algorithmic}
\end{algorithm}

\subsection{棄案：重複執行 hMetis}
在初期的設計中，原本打算利用剩餘的執行時間，加入 Multi-start 的機制重複執行多次完整的 hMetis 流程以尋求最佳解。但經過實驗後發現，最終的 cut size 降低幅度並不明顯，而 runtime 卻直接成倍增長。考量到計分方式中 runtime 佔了 20\% 的權重，為了一點點的 cut size 犧牲大量 runtime 並不划算，因此最終決定主程式只呼叫一次 hMetis。

\section{hMetis 多層次分割流程}

hMetis 結合了 Clustering 與 F-M Partitioning，透過遞迴的方式將龐大的電路圖逐步縮小，求解後再逐步還原微調。

\begin{algorithm}[H]
\caption{hMetis Recursive Flow}
\begin{algorithmic}[1]
\State 執行一次 iteration 的 pre-partition (FM)
\If{pre-partition 失敗 (無法滿足平衡)}
    \State \Return 失敗
\EndIf
\State 執行 First Choice Clustering 將圖縮小
\State 遞迴呼叫 \text{hMetis}
\If{下一層 hMetis 回傳失敗}
    \State 執行 $n$ 次完整的 F-M Partitioning 取最佳解 (Initial Partitioning)
\Else
    \State 將下一層的結果解開 (Declustering)
    \State 執行 F-M Partitioning 進行微調 (Refinement)
\EndIf
\State \Return 本層的 partition 結果
\end{algorithmic}
\end{algorithm}

\subsection{棄案：固定的遞迴停止條件 (Threshold)}
關於 hMetis 的終止條件，一開始是設定當 cell 數量低於一個固定的 threshold 時就停止 clustering 並開始 Initial Partitioning。但後來發現固定 threshold 會造成兩個問題：
\begin{enumerate}
    \item 大型測資在多次 clustering 後，容易出現體積過大的超級節點 (Cluster)，導致後續的 partition 無法找到合乎 balance constraint 的解而失敗。
    \item 小型測資可能因為初始 cell 數量就低於 threshold，導致一次 clustering 都沒做到就直接進入 partition。
\end{enumerate}
基於以上狀況，最後決定讓程式自己動態尋找 hMetis 的深度：持續遞迴進行 clustering，直到下一層的 partition 因為 cluster 太大找不到平衡解而回傳失敗時，才停止遞迴並以此深度作為最底層開始反向 Refinement。

\section{First Choice Clustering}

本階段透過 First Choice 演算法，將連結緊密的 cell 綁定成超級節點，以減少圖的規模並保留自然的群集特徵。主要邏輯與講義概念一致。

\begin{algorithm}[H]
\caption{First Choice Clustering}
\begin{algorithmic}[1]
\For{每一個未被綁定的 cell}
    \State 尋找與其連結的最高 weight 目標 (若平手則選擇 size 較小者)
    \State 將該 cell 併入目標 cluster 中，並更新 cluster size
\EndFor
\State 建立新的 net 與 cluster 映射表，回傳縮小後的圖
\end{algorithmic}
\end{algorithm}

\subsection{略過巨大 Net 的優化}
實作上唯一比較特別的處理是「跳過 cell 數量太多的 net」。因為連接極多 cell 的超級大 net 對於 weight 的計算貢獻值極低 ($1 / (|e|-1)$)，且掃描它們會耗費大量時間，因此在計算權重時直接忽略這些巨型 net，可以大幅加快 clustering 的速度且幾乎不影響最終品質。

\section{F-M Partitioning}

此為最核心的演算法，採用 $O(1)$ 的 Bucket List 資料結構進行節點搬移與 Gain 更新，邏輯同樣基於標準的 F-M 演算法。

\begin{algorithm}[H]
\caption{Fiduccia-Mattheyses Partitioning}
\begin{algorithmic}[1]
\Repeat
    \State 紀錄本回合初始的 cut size 為 \texttt{best\_cut\_size}
    \State 初始化狀態（包含解除所有 cell 的鎖定）
    \State 計算所有 cell 的 initial gain 並建立 Bucket List
    \While{還有未鎖定的 cell 且尚未超時}
        \State 根據策略選出最佳的 cell 進行移動
        \State 將該 cell 移至對立 group 並鎖定
        \State 更新相鄰 cell 的 gain
        \State 檢查並紀錄過程中的 min cut size (需滿足平衡條件)
    \EndWhile
    \State Trace back 回到 min cut size 發生的狀態
\Until{Trace back 後的 cut size $\ge$ \texttt{best\_cut\_size}}
\end{algorithmic}
\end{algorithm}

\subsection{Cell 選擇策略 (Tie-breaking)}
當需要選擇要移動的 cell 時，若兩個 group 的 top cell 都滿足 balance 條件，則依照以下順位進行選擇：
\begin{enumerate}
    \item Gain 較大者優先。
    \item 選擇所在 Group size 較大的一邊，以盡量維持兩邊的平衡。
\end{enumerate}
這樣可以避免未來演算法為了滿足 balance constraint，而被迫選擇 gain 較差的 cell。

\subsection{非法狀態處理}
在 F-M 的過程中，搬移 cell 經常會遇到破壞平衡 (balance constraint) 的情況，這主要基於兩個原因：
\begin{enumerate}
    \item 若 balance degree 設定極小且為偶數個 cell，搬動 1 個 cell 就會違反平衡，必須連續搬動 2 個才能恢復平衡。
    \item 當某個 cell (即 Cluster) 的 size 太大時，移動它雖然能有效降低 gain，但會立刻打破平衡限制。
\end{enumerate}
為了解決這個問題，程式允許這些破壞平衡的 move 發生，但只要移動後的狀態違反 balance，無論當下的 cut size 是否為歷史新低，都不會將其紀錄為 \texttt{min\_cut\_size}。這保證了 iteration 結束後的 trace back 絕對不會回到 invalid 的狀態。若整個 iteration 跑完都沒有出現過任何一次合法的分組，則判定此次 partition 失敗。

\end{document}